\documentclass[12pt]{article}
\usepackage[utf8]{inputenc}
\usepackage[T1]{fontenc}
\usepackage[margin=1in]{geometry}
\usepackage{setspace}
\usepackage{xcolor}
\usepackage[colorlinks=true, urlcolor=blue!60!black, linkcolor=blue!60!black]{hyperref}
\setlength{\parindent}{2em}
\setlength{\parskip}{1.25em}
\renewcommand{\baselinestretch}{1.0}
\begin{document}
\begin{center}
{\LARGE \textit{Troxel v. Granville}, 530 U.S. 57 (2000)\par}
\vspace{0.8em}
{\large Case Brief --- Four-Page Edition\par}
\end{center}
\vspace{1em}

\section{Memory Jogger}

In a fractured 4--2--3 decision, the Court held that a Washington court's application of a ``breathtakingly broad'' nonparental visitation statute to override a fit mother's limits on grandparent visitation violated her fundamental right to direct the care, custody, and control of her children (\href{https://www.law.cornell.edu/supremecourt/text/99-138}{Cornell LII}).

\section{Detailed Case Facts}

Tommie Granville and Brad Troxel never married but had two daughters, Isabelle and Natalie. After they separated, Brad regularly brought the girls to his parents, Jenifer and Gary Troxel. Brad committed suicide in 1993, and Granville later limited the grandparents to one short visit a month and some holidays. The Troxels petitioned under Wash. Rev. Code \S~26.10.160(3), which let ``any person'' petition at any time and let a court order visitation whenever it may serve the child's best interest. They sought two weekends a month and two weeks each summer; Granville offered one day a month without overnights. No one alleged that she was unfit. The trial court ordered one weekend a month, a week each summer, and time on the grandparents' birthdays, presuming visitation to be in the children's best interest and recalling the judge's own childhood visits with his grandparents.

\section{Procedural History}

The Washington Court of Appeals reversed on statutory standing grounds. The Washington Supreme Court affirmed the dismissal 5--4 but held the statute facially unconstitutional, because it required no showing of harm and allowed any person to petition at any time. The U.S. Supreme Court affirmed on June 5, 2000, with no majority opinion.

\section{Judicial Votes}

Affirmed, \textbf{6--3}, across six opinions.
\begin{itemize}
\item \textbf{Plurality:} Justice O'Connor, joined by Chief Justice Rehnquist and Justices Ginsburg and Breyer.
\item \textbf{Concurring in the judgment:} Justice Souter; Justice Thomas.
\item \textbf{Dissenting:} Justices Stevens, Scalia, and Kennedy, each separately.
\end{itemize}

\section{Holding}

Due process protects parents' fundamental right to decide matters of their children's care, custody, and control. Because fit parents are presumed to act in their children's best interests, a court must give at least some special weight to a fit parent's decision. Applying a statute letting any person petition at any time, under a bare best-interest standard, with the presumption inverted against the parent, was unconstitutional as applied. Under \textit{Marks}, that as-applied ruling is the narrowest ground. The Court left open whether a showing of harm is required, the tier of scrutiny, and the validity of narrower statutes.

\section{Analysis of Opinions}

\textbf{Plurality (O'Connor).} Acknowledging grandparents' role in modern families, O'Connor traced parental rights from \textit{Meyer} and \textit{Pierce} through \textit{Yoder}, \textit{Parham}, and \textit{Santosky}. Four factors made the order unconstitutional: Granville was fit; the court gave her judgment no special weight and presumed in favor of the grandparents; it ignored her offer of meaningful visitation; and its slender findings (cousins, music, the judge's memories) showed a mere disagreement about what was better for the girls. The plurality affirmed without remand to spare Granville more litigation.

\textbf{Souter (concurring in the judgment).} The state court's facial invalidation sufficed; the right of upbringing would be a sham if a judge could override a parent whenever he thought he knew better.

\textbf{Thomas (concurring in the judgment).} Taking precedent as given, he would apply strict scrutiny, which Washington could not meet.

\textbf{Stevens (dissenting).} The statute had a plainly legitimate sweep, and the child, not ``so much chattel,'' has interests in established bonds.

\textbf{Scalia (dissenting).} Parental rights are unalienable, but judges may not enforce unenumerated rights against state statutes; legislatures should decide.

\textbf{Kennedy (dissenting).} The harm requirement was wrong; a parent's rights against a stranger differ from those against a de facto parent, and limits should be elaborated with care.

\section{Examples: Future Applications}

\begin{enumerate}
\item \textbf{Same-Side (Invalid): Community Mentors.} Any neighbor or teacher may seek overnight visits whenever a judge finds them enriching, with no deference to fit parents.
\item \textbf{Same-Side (Invalid): Presumption for Grandparents.} A judge orders weekly overnights over a fit widower's objection, presuming grandparents are good for kids unless he proves otherwise.
\item \textbf{Opposite-Side (Valid): Harm Standard.} A statute limited to complete cutoffs after a parent's death, with a presumption for the parent and clear-and-convincing proof of significant harm.
\item \textbf{Opposite-Side (Valid): Unfit Parent.} After a mother is adjudicated unfit, a court orders sibling and grandparent visits; \textit{Troxel}'s presumption protects only fit parents.
\item \textbf{Fence-Sitter: De Facto Parent.} A former partner who co-raised a child for seven years seeks visitation without alleging harm. The plurality reserved the question; Thomas would likely side with the mother, while Stevens and Kennedy flagged de facto parents.
\end{enumerate}

\section{Critique}

\textbf{Originalist views.} Scalia's dissent exemplifies restraint: the Due Process Clause guarantees process, not judge-made family law. Thomas, though skeptical of substantive due process, applied strict scrutiny because no one asked to overrule \textit{Meyer} or \textit{Pierce}.

\textbf{Doctrinal and progressive views.} Emily Buss described parental rights as ``adrift in the middle,'' since the plurality defined neither the standard of review nor special weight, and the 4--2--3 split confused state courts. Woodhouse, Coupet, and Maldonado warn that rigid parental autonomy slights children's bonds in blended and kinship households.

\section{Key Quotations}

\begin{quote}
``The liberty interest at issue in this case---the interest of parents in the care, custody, and control of their children---is perhaps the oldest of the fundamental liberty interests recognized by this Court.''
\end{quote}
\noindent --- Justice O'Connor, plurality opinion.

\begin{quote}
``Accordingly, so long as a parent adequately cares for his or her children (i.e., is fit), there will normally be no reason for the State to inject itself into the private realm of the family to further question the ability of that parent to make the best decisions concerning the rearing of that parent's children.''
\end{quote}
\noindent --- Justice O'Connor, plurality opinion.

\begin{quote}
``As we have explained, the Due Process Clause does not permit a State to infringe on the fundamental right of parents to make childrearing decisions simply because a state judge believes a `better' decision could be made.''
\end{quote}
\noindent --- Justice O'Connor, plurality opinion.

\end{document}
